\documentclass[12pt, ams, nofootinbib]{revtex4}
%\documentclass[12pt,preprint,ams, nofootinbib]{revtex4-1}
%\documentclass{revtex4-1}
\usepackage{amsmath}
\usepackage{amsfonts}
\usepackage{amsthm}
\usepackage{amssymb}
\usepackage{graphicx}  %for \scalebox	
\usepackage{mathtools} %for \dcases
\usepackage{epsfig}
\usepackage[colorlinks=true,urlcolor=blue,linkcolor=blue]{hyperref}
\usepackage{amscd}
%\usepackage[color,draft]{showkeys}
\usepackage{xcolor}
%\usepackage{enumitem}
\usepackage{enumerate}
\usepackage{frcursive}
\usepackage[T1]{fontenc}
\usepackage[toc]{appendix}
%\usepackage[bottom]{footmisc}
\usepackage{mathrsfs}
\usepackage{tensor}     % index positioning in tensor componentes: \tensor{M}{^a_b^c_d}
\usepackage{multirow}

\usepackage{mybookstyle}


\begin{document}

\preprint{APS/123-QED}

\title{Gyroscope precession}

\author{Juan A. Guarín-Rojas}
\email{juan2201870@correo.uis.edu.co}
\affiliation{Escuela de Física, Universidad Industrial de Santander, CP 680002,  Bucaramanga, Colombia.}

\author{Antonio C. Guti\'errez-Pi\~{n}eres}
\email{acgutier@uis.edu.co}
\affiliation{Escuela de Física, Universidad Industrial de Santander, CP 680002,  Bucaramanga, Colombia}
\affiliation{Instituto de Ciencias Nucleares, Universidad Nacional Aut\'onoma de M\'exico,
 AP 70543,  M\'exico, CDMX 04510, M\'exico}

\date{\today}

%%%%%%%%%%%%%%%%%%%%%%%%%%%%%%%%%%%%%%%%%%%%%%%%%%%%%%%%%%
%%%%%%%%%%%%%%%%%%%%%%%%%%%%%%%%%%%%%%%%%%%%%%%%%%%%%%%%%%
%\begin{abstract} 
%\end{abstract}

%\pacs{}

\maketitle

%%%%%%%%%%%%%%%%%%%%%%%%%%%%%%%%%%%%%%%%%%%%%%%%%%%%%%%%%%
%%%%%%%%%%%%%%%%%%%%%%%%%%%%%%%%%%%%%%%%%%%%%%%%%%%%%%%%%%
\section{What is a gyroscope?}

A spinning symmetric top keeps its axis of rotation unchanged with respect to an inertial frame in the absence of torques or in the fast top approximation.

Experimentally, it is convenient to construct a fast top so that, if any torques are present, their effect might be insignificant.

A gimbal has the purpose of controlling or restricting the movement of the top, specifically, restricting the degrees of freedom of the rotation axis.

A gyroscope can be defined as a symmetric top together with a gimbal or a mount that holds the top.



%%%%%%%%%%%%%%%%%%%%%%%%%%%%%%%%%%%%%%%%%%%%%%%%%%%%%%%%%%
%%%%%%%%%%%%%%%%%%%%%%%%%%%%%%%%%%%%%%%%%%%%%%%%%%%%%%%%%%
\section{Gyroscope precession in Newtonian mechanics}
When torques are present, a symmetric top's axis of rotation precesses and nutates around a certain axis. If the angular velocity of the top is large enough, the nutation can be negligible, and the top precesses with a constant angular velocity of $\dot{\phi}=$.

%%%%%%%%%%%%%%%%%%%%%%%%%%%%%%%%%%%%%%%%%%%%%%%%%%%%%%%%%%
%%%%%%%%%%%%%%%%%%%%%%%%%%%%%%%%%%%%%%%%%%%%%%%%%%%%%%%%%%
\section{Gyroscope precession in general relativity}
In general relativity, even in the absence of torques, a gyroscope will precess due to Thompson's precession, geodesic precession, and the Lense-Thirring effect. Thompson's precession manifests the special relativistic nature of spacetime, while the geodesic precession and the Lense-Thirring effect manifest the curvature of spacetime. The Lense-Thirring effect is usually related to the gravitational dragging effect.


%%%%%%%%%%%%%%%%%%%%%%%%%%%%%%%%%%%%%%%%%%%%%%%%%%%%%%%%%%
%%%%%%%%%%%%%%%%%%%%%%%%%%%%%%%%%%%%%%%%%%%%%%%%%%%%%%%%%%
\bibliographystyle{unsrt}
\bibliography{OC_Ref}


\end{document}
